\section{Introduction}

In this paper, we study the modules over face rings, introduced by Athanasiadis and Stanley, whose Hilbert functions are the relative local $h$-vectors of quasi-geometric homology triangulations of simplices, a broad class of formal subdivisions that includes all geometric triangulations and is natural from the point of view of combinatorial commutative algebra.  See Section 2.1 for the precise definition and further references.

Fix an infinite field $k$.  Let $\sigma \colon \Gamma \to 2^V$ be a quasi-geometric homology triangulation of a simplex, and let $E$ be a face of $\Gamma$. Say that a face $G \in \Gamma$ is \emph{interior} if $\sigma(G) = V$, and let $I$ be the ideal in the face ring $k[\lk_\Gamma(E)]$ generated by the faces that are interior relative to $E$, i.e. 
\[
 I = (x^F : F \sqcup E \mbox{ is interior} ).
\]
Let $d = |V|-|E|$, which is the Krull dimension of $k[\lk_\Gamma(E)]$, and let $\theta_1, \ldots, \theta_{d}$ be a special l.s.o.p., as in \cite{Stanley92, Athanasiadis12b}. See also Section~\ref{ss:sr}, where we recall the definition and construction of special l.s.o.p.s.

\begin{definition}
The \emph{local face module} $L(\Gamma,E)$ is defined as the image of $I$ in $k[\lk_\Gamma(E)]/(\theta_1, \ldots, \theta_d)$.
\end{definition}

Note that $L(\Gamma,E)$ is a finite dimensional graded $k$-vector space. The \emph{local $h$-vector} is its Hilbert function: 
\[
\ell(\Gamma, E)  \coloneqq  (\ell_0, \ldots, \ell_{d}), \quad \quad \mbox{ where } \ \ell_i  \coloneqq  \dim L(\Gamma,E)_i.
\]
The local face module $L(\Gamma,E)$ depends on the choice of a special l.s.o.p., but $\ell(\Gamma,E)$ is an invariant of the triangulation with the symmetry $\ell_i = \ell_{d- i}$. See Section~\ref{sec:triangulations} for details and references. In the past few years, there has been significant research activity on the combinatorics of local $h$-vectors and relations to intersection homology \cite{Athanasiadis16, KatzStapledon16, Stapledon17, deCataldoMiglioriniMustata18}.  Recent advances include a proof that every non-negative integer vector satisfying $\ell_0 = 0$ and $\ell_i = \ell_{d-i}$ is the local $h$-vector of a quasi-geometric triangulation for $E = \emptyset$ \cite{JKMS}, and a relative hard Lefschetz theorem that yields unimodality of local $h$-vectors for regular subdivisions in a more general setting (for regular nonsimplicial polyhedral subdivisions that are not necessarily rational) \cite{Karu19}.

 

Here, we investigate the local face modules $L(\Gamma, E)$ using methods of combinatorial commutative algebra. In particular, we describe natural combinatorial resolutions of these modules as well as natural maps of $k[\lk_\Gamma(E)]$-modules, $L(\Gamma,E) \to L(\Gamma, E')$, for $E \subset E'$.  Our first theorem gives explicit generators for the kernel of the natural map $I \to k[\lk_\Gamma(E)]/(\theta_1, \ldots, \theta_d)$. Moreover, we extend this to an exact sequence of graded $k[\lk_\Gamma(E)]$-modules in which each term is a direct sum of degree-shifted monomial ideals.

Label the vertices of the simplex $V = \{ v_1, \ldots, v_n \}$.  For a subset $U \subset V$, let $U^c  \coloneqq  V \smallsetminus U$. After relabeling, we may assume that $\sigma(E)^c = \{v_1, \ldots, v_b\}$. Given $S \subset \{v_1, \ldots, v_d\}$, 
we define the ideal $I_S \subset k[\lk_\Gamma(E)]$ by
\begin{equation*}
I_S  \coloneqq  ( x^F : \, \sigma(F \sqcup E)^c \subset S). 
\end{equation*}
Note that $I_{S'} \subset I_{S}$ for $S' \subset S$, and $I_S$ depends only on $S \cap \{v_1, \ldots, v_b\}$. For instance, $I_{\emptyset} = I$ and $I_{S} = k[\lk_{\Gamma}(E)]$ if $\{v_1, \ldots, v_b \} \subset S$. By the definition of a special l.s.o.p. (Definition~\ref{def:special}), after reordering, we may assume
\[
\supp(\theta_i) \subset \{ w \in \lk_\Gamma(E) : v_i \in \sigma(w) \}, 
\]
for $1 \leq i \leq b$.  As a consequence, for any $v_i \in S$, multiplication by $\theta_i$ induces a degree 1 map $\lambda_i \colon I_S \to I_{S \smallsetminus \{v_i\}}$.

\begin{theorem}\label{thm:resolution}
There is an exact sequence of graded $k[\lk_{\Gamma}(E)]$-modules
 $$0  \to k[\lk_{\Gamma}(E)][-d] \to \bigoplus_{ |S|  = d - 1 } I_{S}[-(d-1)] \to \dotsb  \to \bigoplus_{|S| = 1} I_{S}[-1] \to I \to L(\Gamma,E) \to 0,$$
where,  
for $S= \{v_{i_0}, \ldots, v_{i_k}\}$, with $i_0 < \cdots < i_k$, the differential restricted to $I_S$ is $\oplus_{j = 0}^k (-1)^j \lambda_{i_j}$.
\end{theorem}

\begin{corollary}\label{cor:presentation}
The kernel of the surjection $I \to L(\Gamma, E)$ is the ideal $J$ generated by 
 $$ \left \{ \theta_i \cdot x^{F} : F \sqcup E \mbox{ is interior } \right \}  \cup  \left \{ \theta_{j} \cdot x^G : \sigma(G \sqcup E) = \{v_j\}^c, \mbox{ for } 1 \leq j \leq b \right \}.$$
\end{corollary}

We also construct maps between local face modules, as follows. For faces $E \subset E'$ in $\Gamma$, let $\Star(E' \smallsetminus E)$ denote the closed star of $E' \smallsetminus E$ in $\lk_\Gamma(E)$. 
We have a natural inclusion of complexes 
$\lk_\Gamma(E') \subset \lk_\Gamma(E)$. 




\begin{theorem}\label{thm:maps}
Let $E \subset E'$ be faces of $\Gamma$, with $$d = n - \vert E \vert, \quad  
d' = n - \vert E' \vert, \quad 
\mbox{ and } \quad b' = n - |\sigma(E')|.$$ Let $\{\theta_1, \dotsc, \theta_{d } \}$ be a special l.s.o.p. for $k[\lk_{\Gamma}(E)]$,  
and let 
$\theta_i'  \coloneqq  \theta_i|_{\Star(E' \smallsetminus E)}$.  Then there is a 
unique homomorphism of graded 
$k$-algebras 
\[
\phi\colon k[\lk_{\Gamma}(E)]/(\theta_1, \dotsc, \theta_d) \to k[\lk_{\Gamma}(E')]/(k[\lk_\Gamma(E')]  \cap (\theta_1', \dotsc, \theta_{d}'))
\]
whose kernel contains $\{[x^F] : F \not \in \Star(E' \smallsetminus E)\}$ and satisfies $\phi(x^F) = x^F$ for all $F \in \lk_{\Gamma}(E')$. Moreover, there is a special l.s.o.p. $\zeta_1, \ldots, \zeta_{d'}$ for $k[\lk_\Gamma(E')]$ such that $(\zeta_1, \ldots, \zeta_{d'}) = k[\lk_\Gamma(E')] \cap (\theta'_1, \ldots, \theta'_{d})$ and, up to reordering, we have $\theta_i|_{\lk_\Gamma(E')} = \zeta_i$, for $1 \leq i \leq b'$.  With this choice of special l.s.o.p., $\phi(L(\Gamma,E)) \subset L(\Gamma,E')$. 
\end{theorem}

\begin{remark}
Theorem~\ref{thm:maps} may be viewed as a functoriality statement for local face modules. Start by fixing the special l.s.o.p. $\theta_1, \ldots, \theta_d$.  Then $L(\Gamma, E)$ is well-defined. For $E' \supset E$ the special l.s.o.p. $\zeta_1, \ldots, \zeta_{d'}$ depends on some choices, but the ideal that it generates does not, nor does the map $\phi \colon L(\Gamma,E) \to L(\Gamma,E')$.  Moreover, for $E'' \supset E'$, one readily checks that the maps $\phi' \colon L(\Gamma, E') \to L(\Gamma, E'')$ and $\phi '' \colon L(\Gamma,E) \to L(\Gamma, E'')$ are independent of all choices and satisfy $\phi'' = \phi' \circ \phi$.  Thus one obtains a functor from the poset of faces of $\Gamma$ that contain $E$ to graded vector spaces, given by $E' \mapsto L(\Gamma, E')$.
\end{remark}

We now give two applications of the above theorems. The first is a monotonicity property for local $h$-vectors.

\begin{theorem}\label{thm:increase} 
Let $E \subset E'$ be faces of $\Gamma$ such that $\sigma(E) = \sigma(E')$. Then $\ell (\Gamma, E) \geq \ell(\Gamma, E')$.
\end{theorem}

The inequality in Theorem~\ref{thm:increase} is term by term, i.e. $\dim L(\Gamma, E)_i \geq \dim L(\Gamma,E')_i$ for all $i$. The proof is by showing that the map $\phi \colon L(\Gamma, E) \to L(\Gamma,E')$ given by Theorem~\ref{thm:maps} is surjective.


Our second application of the above theorems is to a decades old problem posed by Stanley, who introduced and studied local $h$-vectors in the special case where $E = \emptyset$ and asked for a characterization of triangulations for which they vanish \cite[Problem~4.13]{Stanley92}. This problem remains open, and is of enduring interest \cite[Problem~2.12]{Athanasiadis16}. The extension to the case where $E$ is not empty is particularly relevant for applications to the monodromy conjecture \cite{Igusa78, DenefLoeser98, Stapledon17}. In \cite{LarsonPayneStapledon}, we prove a theorem on the structure of geometric triangulations with vanishing local $h$-vectors that is tailored to this purpose, and we use it to prove the monodromy conjectures for all singularities that are nondegenerate with respect to a simplicial Newton polyhedron.  See Theorems~1.1.1, 1.4.5, and 4.1.3 in loc. cit..

Here, we apply  Theorem~\ref{thm:resolution} to prove another theorem on the structure of faces in triangulations with vanishing local $h$-vectors.  Let $F \in \lk_\Gamma(E)$ be a face such that $F\sqcup E$ is interior. Following terminology from the monodromy conjecture literature (see, e.g. \cite{LemahieuVanProeyen11}), we say that $F$ is a \emph{pyramid with apex $w \in F$}  if $(F \sqcup E )  \smallsetminus w$ is not interior. Let $$\mathcal{A}_F  \coloneqq  \{ w \in F : F \mbox{ is a pyramid with apex } w \}, \mbox{ \ and \ } V_w  \coloneqq  \sigma( (F \sqcup E) \smallsetminus w)^c.$$
The elements of $V_w$ correspond to the \emph{base directions} of $F$, i.e. the facets of $2^V$ that contain the base of $F$, when viewed as a pyramid with apex $w$.  We say $F$ is a $U$-pyramid if there is an apex $w \in \mathcal{A}_F$ such that $|V_w|  = 1$. In other words, a $U$-pyramid is a pyramid with a unique base direction, for some choice of apex.

\begin{definition}
Let $F \in \lk_{\Gamma}(E)$ be a face.  An \emph{interior partition} of $F$ is a decomposition
\[
F = F_1 \sqcup F_2 \sqcup \mathcal{A}_F
\]
such that $F_1 \sqcup \mathcal{A}_F \sqcup E$ and $F_2 \sqcup \mathcal{A}_F \sqcup E$ are both interior.
\end{definition}

\begin{theorem}\label{thm:interiornonvanish}
Suppose $\ell(\Gamma,E) = 0$ and $F \in \lk_\Gamma(E)$ has an interior partition $F = F_1 \sqcup F_2 \sqcup \mathcal{A}_F$ such that $|F_1| \leq 2$.  Then $F$ is a $U$-pyramid.
\end{theorem}

See Remark~\ref{r:specialcase} for a short proof in a special case that illustrates the naturality of the 
$U$-pyramid condition.
The method of proof breaks down when $|F_i| \geq 3$.  See Example~\ref{example:nonrestriction}.  


\begin{remark}
The analogous theorem in \cite{LarsonPayneStapledon} requires that the triangulation be geometric and that the interior partition satisfies the additional condition $\sigma(F_2 \sqcup E)^c = \bigcup_{w \in \mathcal{A}_F} V_w$.  But then the hypothesis that $|F_1| \leq 2$ is dropped entirely.  So, even for geometric triangulations, there are cases of Theorem~\ref{thm:interiornonvanish} that are not necessarily covered by \cite[Theorem~4.1.3]{LarsonPayneStapledon}. It should be interesting to look for a common generalization of these vanishing results, and to pursue further progress on Stanley's problem of characterizing triangulations with vanishing local $h$-vector more generally.
\end{remark}


\begin{remark}
To the best of our knowledge, all of the theorems stated in the introduction are new even for regular triangulations. The reader who prefers to do so may safely restrict attention to geometric or even regular triangulations.  However, while the structure results for triangulations with vanishing local $h$-vectors in \cite{dMGPSS20} and \cite{LarsonPayneStapledon} rely on special properties of geometric triangulations, the proofs presented here work equally well for quasi-geometric homology triangulations, and we find it  natural to work in this level of generality.
\end{remark}
